% From_Bit_To_Boundary_V4.tex
% Reconstructed from the typeset PDF (From_Bit_to_Boundary_V4-1.pdf, v4.0, April 2026).
% The original LaTeX source was not available. Text, equations, numbering, and
% references follow the PDF; layout details may differ. Proofread against the PDF.

\documentclass[11pt]{article}

\usepackage[margin=1in]{geometry}
\usepackage[T1]{fontenc}
\usepackage[utf8]{inputenc}
\usepackage{amsmath, amssymb, amsthm}
\usepackage{booktabs}
\usepackage{array}
\usepackage{hyperref}

\hypersetup{
  colorlinks=true,
  linkcolor=blue,
  citecolor=blue,
  urlcolor=blue
}

\theoremstyle{plain}
\newtheorem{assumption}{Assumption}
\newtheorem{remark}{Remark}[section]
\newtheorem{definition}{Definition}[section]
\newtheorem{lemma}{Lemma}[section]
\newtheorem{proposition}{Proposition}[section]
\newtheorem{corollary}{Corollary}[section]
\newtheorem{theorem}{Theorem}[section]

\DeclareMathOperator*{\argmax}{arg\,max}

\emergencystretch=2em

\title{From Bit to Boundary\\[4pt]
  \large A Geometric Theory of Information Collapse}
\author{Amber Anson\textsuperscript{1} \quad Claude\textsuperscript{2}\\[2pt]
  \small \textsuperscript{1}AmberContinuum Research \quad \textsuperscript{2}Anthropic\\
  \small \texttt{ambercontinuum@gmail.com}}
\date{April 2026 --- v4.0}

\begin{document}
\maketitle

\begin{abstract}
Shannon's information theory (1948) formalized transmission: entropy, channel capacity,
and coding limits. By design, Shannon excluded semantics---how transmitted signals acquire
definite meaning. This paper introduces a complementary geometric principle---the Post-Shannon
Law---governing the conditions under which latent information collapses into stable, observable
outcomes.

We define a Universal Collapse Operator $\Phi(M, T, R, B, \varepsilon, \tau) \to S$ over a compact
Riemannian state manifold and derive the collapse boundary condition $R - \tau \geq \varepsilon$,
where $R$ is resonance coupling (a strictly log-concave kernel), $\tau$ is torsion (internal
contradiction cost, modeled as a filtered Poisson process), and $\varepsilon$ is the collapse
threshold. We prove uniqueness, irreversibility, and torsion monotonicity under explicit regularity
assumptions. The operator exhibits structural isomorphism across cognitive, computational,
historical, and psychological substrates.

We situate this framework within the cybernetic tradition of Wiener, Ashby, and Bateson,
identifying the collapse boundary as a generalization of homeostatic regulation: where cybernetics
describes the maintenance of viable states through feedback, collapse dynamics describe the
irreversible transition from superposed latent information to determinate structure. Substrate
independence is evidenced by structural isomorphism with quantum measurement, perceptual
bistability, bounded-confidence opinion dynamics, and belief revision---each an independent line
of published research in which the six-component collapse structure is already implicit. Testable
predictions and a pre-registered empirical protocol are specified.
\end{abstract}

\noindent\textbf{Keywords:} information theory, information collapse, cybernetics, torsion,
geometric boundaries, substrate independence, semantic emergence, requisite variety, homeostatic
regulation

% ============================================================
\section{Introduction}
% ============================================================

\subsection{The Shannon Formalism and Its Scope}

Shannon's ``A Mathematical Theory of Communication'' (1948) solved transmission: given a source
of messages, a noisy channel, and an encoder--decoder pair, Shannon established entropy $H(X)$ as
the measure of source uncertainty, channel capacity $C$ as the supremum of reliable transmission
rates, and the source and channel coding theorems as fundamental limits~\cite{shannon1948}. These
results are substrate-agnostic in the transmission domain---they hold for electrical signals,
optical fibers, and biological channels alike.

Shannon was explicit about what his formalism excluded:
\begin{quote}
``Frequently the messages have meaning; that is they refer to or are correlated according
to some system with certain physical or conceptual entities. These semantic aspects of
communication are irrelevant to the engineering problem.''~\cite{shannon1948}
\end{quote}

This exclusion was methodologically sound---separating transmission fidelity from interpretation
permitted rigorous treatment of the former. But it left a structural gap: the formalism provides
no account of how transmitted information acquires definite semantic content, how superposed
interpretations resolve, or what governs the transition from probabilistic latent states to
determinate outcomes.

\subsection{Cybernetics and the Regulation Problem}

The cybernetic tradition addressed a related but distinct problem: how systems maintain viability
under perturbation. Wiener's (1948) formalization of feedback---circular causality between
controller and controlled---established that stable behavior requires information about outcomes
to flow back and modify inputs~\cite{wiener1948}. Ashby's Law of Requisite Variety (1956) proved
that a regulator must possess at least as many distinguishable states as the disturbances it must
absorb: only variety can absorb variety~\cite{ashby1956}. Bateson (1972) extended cybernetic
reasoning into epistemology, arguing that ``the difference which makes a difference'' is the
elementary unit of information in any living system~\cite{bateson1972}.

These contributions share a structural feature: they describe \emph{maintenance of viable states
within boundaries}. Homeostasis is the persistence of a system's essential variables within limits
compatible with continued operation. The cybernetic framework accounts for how systems remain
in viable regions of state space, but does not formalize the irreversible transition from latent
superposition to collapsed structure---the moment when maintenance fails or is unnecessary because
a new determinate state has crystallized.

\subsection{The Present Contribution}

This paper formalizes \emph{information collapse}: the geometric conditions under which superposed
latent information transitions irreversibly into a single determinate outcome. The collapse
boundary $R - \tau \geq \varepsilon$ extends cybernetic boundary logic from maintenance to phase
transition. Where Ashby's variety constraint describes what a regulator must possess to maintain
homeostasis, the collapse threshold $\varepsilon$ describes what must be exceeded for latent
information to become determinate structure. Torsion $\tau$, the internal contradiction cost,
functions as the information-geometric analog of Ashby's disturbance variety---unresolved
contradiction that must be overcome for collapse to proceed.

The framework is instantiated across four independent domains---quantum measurement, perceptual
bistability, opinion dynamics, and belief revision---each drawn from published research by other
investigators in which the collapse operator structure is already implicit
(Section~\ref{sec:substrate}). A pre-registered empirical protocol for the AI safety application is
specified under CDOT 2.0 (osf.io/7mejx).

% ============================================================
\section{The Collapse Operator: Formal Specification}
% ============================================================

\subsection{Standing Assumptions}

We begin by stating the regularity conditions under which all subsequent results hold.

\begin{assumption}[Manifold regularity]\label{asm:manifold}
$M$ is a compact, connected, smooth ($C^\infty$) Riemannian manifold of dimension $n$, equipped
with the metric $g$. Compactness guarantees that continuous functions on $M$ attain their extrema
and that all relevant integrals converge. We further assume $M$ has non-negative sectional
curvature ($\sec_M \geq 0$), which is required for the log-concavity lemma
(Lemma~\ref{lem:logconcave}).
\end{assumption}

\begin{remark}[Relaxation of compactness and curvature]
Compactness can be relaxed to $\sigma$-compactness if $R$ and $\tau$ decay sufficiently at
infinity. In the quantum measurement instantiation (Section~\ref{sec:quantum}), $M$ is a Hilbert
space (non-compact), but the relevant dynamics are confined to a finite-dimensional subspace spanned
by the measurement basis, which is compact. For curvature, the non-negativity condition can be
relaxed: on a Cartan--Hadamard manifold (simply connected, sectional curvature $\leq 0$), the
squared geodesic distance from any point is a globally geodesically convex
function~\cite{petersen2016}, so the log-concavity lemma holds without any segment-length
restriction. On manifolds with sectional curvature bounded above by $K > 0$, geodesic convexity of
the squared distance holds on geodesic balls of radius less than $\pi/(2\sqrt{K})$ (Rauch
comparison).
\end{remark}

\begin{assumption}[Transformation regularity]\label{asm:transformation}
$T : M \to M$ is a $C^1$ diffeomorphism preserving a smooth measure $\mu$ on $M$. This ensures
that the state evolution $\psi(t+1) = T(\psi(t))$ neither creates nor destroys probability mass.
Periodicity of $T$ is not required but is satisfied in several instantiations (e.g., Hamiltonian
evolution in quantum mechanics).
\end{assumption}

\begin{assumption}[Basis genericity]\label{asm:genericity}
The symbolic basis $B = \{b_1, \ldots, b_k\} \subset \mathbb{R}^m$ is in general position: no two
basis elements are equidistant from any point in the image of $M$ under the embedding
$\iota : M \hookrightarrow \mathbb{R}^m$. Formally, for all $\psi \in M$ and all $i \neq j$,
\[
  \|\iota(\psi) - b_i\| \neq \|\iota(\psi) - b_j\|.
\]
This is a generic condition: the set of bases violating it has Lebesgue measure zero in
$(\mathbb{R}^m)^k$. In the quantum measurement case, general position corresponds to the
non-degeneracy of the observable's eigenvalue spectrum.
\end{assumption}

\begin{assumption}[Resonance regularity]\label{asm:resonance}
The resonance function $R : M \times B \to [0, \infty)$ is:
\begin{enumerate}
  \item \emph{Bounded:} $\sup_{\psi \in M, b \in B} R(\psi, b) < \infty$ (guaranteed by
    compactness of $M$ and finiteness of $B$).
  \item \emph{Strictly log-concave in $\psi$:} for fixed $b$, the map $\psi \mapsto \log R(\psi, b)$
    is strictly concave along geodesics in $M$. This is satisfied by Gaussian kernels on
    non-negatively curved manifolds (Lemma~\ref{lem:logconcave}) and by the Born rule
    $R = |c_j|^2$ in finite-dimensional Hilbert spaces.
  \item \emph{Continuous:} $R$ is jointly continuous in both arguments.
\end{enumerate}
\end{assumption}

\begin{assumption}[Torsion regularity]\label{asm:torsion}
$\tau : M \to \mathbb{R}^+$ is continuous and bounded: $0 \leq \tau(\psi) \leq \tau_{\max}$ for all
$\psi \in M$. Boundedness follows from the Poisson torsion model (Section~\ref{sec:poisson}) when
the contradiction arrival rate $\alpha$ is finite and the regularization $\lambda > 0$.
\end{assumption}

\begin{remark}
Assumptions~\ref{asm:manifold}--\ref{asm:torsion} are sufficient for all theorems in
Section~\ref{sec:theorems}. Domain-specific instantiations (Section~\ref{sec:substrate}) may
satisfy stronger conditions; we state only what is needed.
\end{remark}

\subsection{Definitions}

\begin{definition}[Universal Collapse Operator]
Under Assumptions~\ref{asm:manifold}--\ref{asm:torsion}, information collapse is governed by the
operator
\[
  \Phi : (M, T, R, B, \varepsilon, \tau) \to S
\]
where $M$ is the state manifold, $T$ is the transformation, $R$ is the resonance coupling, $B$ is
the symbolic basis, $\varepsilon \in \mathbb{R}^+$ is the collapse threshold, $\tau$ is the torsion
function, and $S \in B$ is the collapsed state.
\end{definition}

\begin{definition}[Collapse Necessity]
Stable state $S$ emerges from latent manifold $M$ if and only if
\[
  \exists\, b \in B \text{ such that } R(\psi, b) - \tau(\psi) \geq \varepsilon.
\]
\end{definition}

\begin{definition}[Torsion]
Torsion $\tau$ quantifies the energetic cost of maintaining contradictory information states:
\begin{equation}\label{eq:torsion}
  \tau(\psi) = \sum_{i=1}^{N_c(\psi)} \frac{1}{(\Delta t_i + \lambda)^2},
\end{equation}
where $N_c(\psi)$ is the number of active contradictions in state $\psi$, $\Delta t_i$ is the
temporal distance between the $i$-th contradictory pair, and $\lambda > 0$ is a regularization
constant preventing divergence at $\Delta t_i = 0$. In differential geometry, torsion measures how
a connection twists parallel-transported vectors off the tangent plane. In information geometry, it
measures the internal tension of self-contradictory states---the structural barrier preventing
coherent collapse.
\end{definition}

\begin{definition}[Collapse Projection]
Define the collapse map $P_B : M \to B$ by
\[
  P_B(\psi) = \argmax_{b \in B} R(\psi, b).
\]
Under Assumptions~\ref{asm:resonance} and the basis genericity assumption, $P_B$ is well-defined
(Lemma~\ref{lem:argmax}). The collapse map induces an equivalence relation $\sim_B$ on $M$:
$\psi_1 \sim_B \psi_2$ iff $P_B(\psi_1) = P_B(\psi_2)$. The quotient space $M/{\sim_B}$ has exactly
$|B| = k$ equivalence classes, each a basin of collapse for the corresponding basis element.
\end{definition}

\begin{definition}[Substrate Independence]
The collapse operator $\Phi$ exhibits structural isomorphism across substrates: there exist
domain-specific encodings $\iota_d : M_d \hookrightarrow \mathbb{R}^{n_d}$ for each domain $d$ such
that the six-component structure $(M_d, T_d, R_d, B_d, \varepsilon_d, \tau_d)$ and the boundary
condition $R_d - \tau_d \geq \varepsilon_d$ are preserved under the encoding.
\end{definition}

\subsection{Resonance Kernel}

The specific resonance function used throughout this paper is the Gaussian kernel on $M$:
\begin{equation}\label{eq:resonance}
  R(\psi, b) = \exp\!\left( -\frac{d_g(\psi, \pi_b(\psi))^2}{2\sigma^2} \right),
\end{equation}
where $d_g$ is the geodesic distance on $M$, $\pi_b(\psi)$ is the geometric projection of $\psi$
onto basis element $b$, and $\sigma > 0$ is the bandwidth parameter. A phase-alignment factor
$\cos(\Delta\phi(\psi, b))$ may be included in domain-specific instantiations where coherent
oscillatory structure is present (e.g., quantum measurement); its inclusion requires an additional
regularity condition on the phase map that is beyond the scope of the substrate-independent
formalism and we do not rely on it in what follows.

\begin{lemma}[Strict log-concavity of the Gaussian kernel]\label{lem:logconcave}
The map $\psi \mapsto R(\psi, b) = \exp(-d_g(\psi,\allowbreak \pi_b(\psi))^2 /\allowbreak 2\sigma^2)$ is strictly
log-concave along geodesics in $M$ whenever the squared distance $\psi \mapsto d_g(\psi, q)^2$ is
strictly geodesically convex for every fixed point $q \in M$ on the geodesic segments under
consideration.
\end{lemma}

\begin{proof}
Fix $b$ and write $q = \pi_b(\psi)$. Then
\[
  \log R(\psi, b) = -\frac{d_g(\psi, q)^2}{2\sigma^2}.
\]
This is the composition of $\psi \mapsto d_g(\psi, q)^2$ with multiplication by the negative
constant $-1/(2\sigma^2)$. Strict geodesic convexity of the squared distance therefore yields strict
geodesic concavity of $\log R$, which is the required log-concavity of $R$. The squared-distance
function is known to be strictly geodesically convex on Cartan--Hadamard manifolds globally, and on
non-negatively curved manifolds along geodesic segments on which no conjugate point is crossed;
see~\cite[Chapter~6]{petersen2016}. In each of these cases the lemma's hypothesis is satisfied.
\end{proof}

\subsection{Torsion as a Filtered Poisson Process}\label{sec:poisson}

The phenomenological equation $d\tau/dt = \alpha' - \beta\tau$ arises as a mean-field approximation
to a stochastic contradiction process. We state the model, derive the stationary expectation and
variance in closed form, and then derive the approximation explicitly with its regime of validity.

\begin{definition}[Contradiction process]
Let $\{N(t)\}_{t \geq 0}$ be a homogeneous Poisson process with rate $\alpha > 0$, representing the
arrival of new contradictions. Each contradiction $i$ arriving at time $t_i$ contributes a decaying
cost
\[
  c_i(t) = \frac{1}{((t - t_i) + \lambda)^2}, \qquad t \geq t_i,
\]
where $\lambda > 0$ is a regularization constant ensuring $c_i(t_i) = 1/\lambda^2 < \infty$. The
instantaneous torsion is the superposition of all active contradiction costs:
\begin{equation}\label{eq:shotnoise}
  \tau(t) = \sum_{i : t_i \leq t} c_i(t) = \sum_{i : t_i \leq t} \frac{1}{(t - t_i + \lambda)^2}.
\end{equation}
This is a shot-noise process~\cite{daley2003}.
\end{definition}

\begin{proposition}[Stationary torsion moments]\label{prop:moments}
As $t \to \infty$, the expected torsion and torsion variance converge to:
\begin{align}
  \mathbb{E}[\tau_\infty] &= \frac{\alpha}{\lambda}, \label{eq:mean} \\
  \operatorname{Var}[\tau_\infty] &= \frac{\alpha}{3\lambda^3}. \label{eq:var}
\end{align}
\end{proposition}

\begin{proof}
We apply Campbell's theorem~\cite{kingman1993}: for a homogeneous Poisson process with rate
$\alpha$ and response function $h(s) = (s + \lambda)^{-2}$, the shot-noise process
$\tau(t) = \sum_{i : t_i \leq t} h(t - t_i)$ has
\[
  \mathbb{E}[\tau(t)] = \alpha \int_0^t h(s)\, ds, \qquad
  \operatorname{Var}[\tau(t)] = \alpha \int_0^t h(s)^2\, ds.
\]
For the mean:
\[
  \int_0^t \frac{ds}{(s + \lambda)^2}
  = \left[ -\frac{1}{s + \lambda} \right]_0^t
  = \frac{1}{\lambda} - \frac{1}{t + \lambda}
  \xrightarrow{t \to \infty} \frac{1}{\lambda}.
\]
For the variance:
\[
  \int_0^t \frac{ds}{(s + \lambda)^4}
  = \left[ -\frac{1}{3(s + \lambda)^3} \right]_0^t
  = \frac{1}{3\lambda^3} - \frac{1}{3(t + \lambda)^3}
  \xrightarrow{t \to \infty} \frac{1}{3\lambda^3}.
\]
Multiplying by $\alpha$ yields Eqs.~\eqref{eq:mean}--\eqref{eq:var}.
\end{proof}

\begin{corollary}[Mean-field approximation]\label{cor:meanfield}
The expected torsion $\bar\tau(t) = \mathbb{E}[\tau(t)]$ satisfies
\begin{equation}\label{eq:exactode}
  \frac{d\bar\tau}{dt} = \alpha\, h(t) = \frac{\alpha}{(t + \lambda)^2}.
\end{equation}
The filtered Poisson process $\tau(t)$ is not Markovian in $\tau$ alone (its future depends on the
full arrival history, not only the current value), so the first-order ODE
\begin{equation}\label{eq:meanfield}
  \frac{d\bar\tau}{dt} \approx \alpha' - \beta\, \bar\tau, \qquad
  \alpha' = \frac{\alpha}{\lambda^2}, \quad \beta = \frac{1}{\lambda},
\end{equation}
is a mean-field approximation that matches the stationary moments
(Proposition~\ref{prop:moments}) and the initial rate
$d\bar\tau/dt|_{t=0} = \alpha/\lambda^2$, but does not reproduce the exact transient dynamics of
Eq.~\eqref{eq:exactode}. The approximation is closest for $t \gg \lambda$, where both
Eq.~\eqref{eq:exactode} and Eq.~\eqref{eq:meanfield} give $\bar\tau \to \alpha/\lambda$, and for
regimes where the response function's memory kernel is short relative to the observation timescale.
\end{corollary}

\begin{proof}
Differentiating $\bar\tau(t) = \alpha \int_0^t h(s)\, ds$ gives $d\bar\tau/dt = \alpha h(t)$,
establishing Eq.~\eqref{eq:exactode}. The mean-field form (Eq.~\eqref{eq:meanfield}) is obtained by
substituting the stationary mean $\bar\tau_{\mathrm{eq}} = \alpha/\lambda$ as the fixed point and
choosing $\beta = 1/\lambda$ so that the linearized relaxation rate matches the scaling set by
$\lambda$. The regime-of-validity statement follows by comparing the two ODEs at short times
($t \ll \lambda$: the exact ODE gives $d\bar\tau/dt \approx \alpha/\lambda^2$, matching the
approximation at $\bar\tau = 0$) and long times (both yield the same stationary value).
\end{proof}

\begin{remark}
The variance result (Eq.~\eqref{eq:var}) provides a principled confidence interval for torsion
estimates: $\tau \in \alpha/\lambda \pm 1.96\sqrt{\alpha/3\lambda^3}$ at 95\% coverage. This
predicts that instruments measuring torsion in systems near the collapse boundary should exhibit
measurable classification variance of order $\sqrt{\alpha/3\lambda^3}$---a feature, not a defect,
of the stochastic model.
\end{remark}

\subsection{Key Theorems}\label{sec:theorems}

\begin{lemma}[Uniqueness of arg max]\label{lem:argmax}
Under Assumption~\ref{asm:resonance} (strict log-concavity) and the basis genericity assumption,
for any $\psi \in M$ there exists a unique $b^* \in B$ such that $R(\psi, b^*) > R(\psi, b)$ for all
$b \in B \setminus \{b^*\}$.
\end{lemma}

\begin{proof}
Suppose for contradiction that $R(\psi, b_i) = R(\psi, b_j)$ for some $i \neq j$. The
equal-resonance set
$\{\psi \in M : R(\psi, b_i) = R(\psi, b_j)\} = \{\psi : \log R(\psi, b_i) - \log R(\psi, b_j) = 0\}$
is the zero set of the smooth function $\psi \mapsto \log R(\psi, b_i) - \log R(\psi, b_j)$. By
strict log-concavity (Assumption~\ref{asm:resonance}) and general position (no two basis elements
are equidistant from any $\psi \in M$ in the induced geometry), this zero set is a smooth
hypersurface of codimension one, which has $\mu$-measure zero. The generic-position condition
ensures that $\psi$ does not lie on this hypersurface, yielding a contradiction. Hence
$b^* = \argmax_{b \in B} R(\psi, b)$ is unique for all $\psi \in M$.
\end{proof}

\begin{theorem}[Uniqueness of Collapse]\label{thm:uniqueness}
Under Assumptions~\ref{asm:manifold}--\ref{asm:torsion}, if collapse occurs (i.e.,
$\exists\, b \in B$ with $R(\psi, b) - \tau(\psi) \geq \varepsilon$), the outcome $S = P_B(\psi)$
is unique.
\end{theorem}

\begin{proof}
By Lemma~\ref{lem:argmax}, $b^* = \argmax_{b \in B} R(\psi, b)$ is unique. If collapse occurs,
$R(\psi, b^*) - \tau(\psi) \geq \varepsilon$. Since $b^*$ is the unique maximizer,
$R(\psi, b^*) > R(\psi, b)$ for all $b \neq b^*$, so $S = b^*$ is the unique collapse outcome.
\end{proof}

\begin{theorem}[Irreversibility]\label{thm:irreversibility}
The collapse map $P_B : M \to B$ is non-injective. Each equivalence class
$[b_j] = P_B^{-1}(b_j)$ is the interior of a generalized Voronoi cell under the resonance-weighted
metric; each such cell is open in $M$, has positive $\mu$-measure, and has topological dimension
$n$.
\end{theorem}

\begin{proof}
Define $V_j = \{\psi \in M : R(\psi, b_j) > R(\psi, b_i) \text{ for all } i \neq j\}$. By
Lemma~\ref{lem:argmax}, the sets $V_j$ are precisely the fibers $P_B^{-1}(b_j)$ (excluding the
measure-zero set of ties, which is empty under general position). Each $V_j$ is open: if
$\psi_0 \in V_j$, then $R(\psi_0, b_j) - R(\psi_0, b_i) > 0$ for all $i \neq j$; by continuity of
$R$ in $\psi$ (Assumption~\ref{asm:resonance}), there is an open neighborhood of $\psi_0$ on which
this inequality is preserved, and this neighborhood lies in $V_j$. Openness in $M$ implies
$\dim V_j = n$ and $\mu(V_j) > 0$ whenever $V_j$ is non-empty (which holds for each $j$ by the
general-position assumption and the fact that $R(\cdot, b_j)$ achieves its maximum on $M$ by
compactness). Since $\dim(M) = n$ and $|B| = k < \infty$, the map $P_B$ sends an $n$-dimensional
space to a finite set; non-injectivity is immediate. Information about which $\psi \in V_j$
obtained prior to collapse is irretrievable from knowledge of $S = b_j$ alone.
\end{proof}

\begin{theorem}[Torsion Monotonicity]\label{thm:monotonicity}
Let $\psi$ be drawn from a probability distribution on $M$ with smooth positive density $p(\psi)$
with respect to $\mu$. For fixed $\varepsilon > 0$, consider uniform shifts of the torsion
function: $\tau(\psi) \mapsto \tau(\psi) + \delta$ for $\delta > 0$. The collapse probability
\[
  P(\delta) = \int_M \mathbf{1}\!\left[ \max_{b \in B} R(\psi, b) - \tau(\psi) - \delta \geq
  \varepsilon \right] p(\psi)\, d\mu(\psi)
\]
is strictly decreasing in $\delta$.
\end{theorem}

\begin{proof}
Define the collapse set
$C(\delta) = \{\psi \in M : \max_b R(\psi, b) - \tau(\psi) - \delta \geq \varepsilon\}$. For
$\delta_1 < \delta_2$, the defining inequality is strictly more restrictive at $\delta_2$, so
$C(\delta_2) \subseteq C(\delta_1)$. The set-difference
\[
  A = C(\delta_1) \setminus C(\delta_2)
  = \{\psi : \varepsilon + \delta_1 \leq \max_b R(\psi, b) - \tau(\psi) < \varepsilon + \delta_2\}
\]
is an open set in $M$: it is the preimage under the continuous function
$\psi \mapsto \max_b R(\psi, b) - \tau(\psi)$ of the half-open interval
$[\varepsilon + \delta_1, \varepsilon + \delta_2)$, whose interior is
$(\varepsilon + \delta_1, \varepsilon + \delta_2)$; the preimage of this interior is open, and by
continuity is non-empty whenever the function's range includes the interval, which holds on any
$M$ not entirely in a single collapse regime. Since $p$ is strictly positive, $\mu(A) > 0$ implies
$\int_A p\, d\mu > 0$, hence $P(\delta_1) - P(\delta_2) = \int_A p\, d\mu > 0$.
\end{proof}

% ============================================================
\section{Torsion as Anti-Bit}
% ============================================================

Shannon's bit is the elementary unit of uncertainty reduction: one bit resolves one binary
distinction, decreasing entropy by $\Delta H = -1$. We define torsion as its structural complement
in the collapse domain: one unit of torsion introduces one unit of internal contradiction,
increasing the cost of collapse by $\Delta\tau = +1$. Where bits enable transmission, torsion
prevents collapse.

\begin{table}[h]
\centering
\caption{Bit--Anti-Bit duality across transmission and collapse domains.}
\begin{tabular}{lll}
\toprule
\textbf{Property} & \textbf{Bit (Shannon)} & \textbf{Anti-Bit / Torsion (Post-Shannon)} \\
\midrule
Domain            & Transmission              & Collapse \\
Effect on entropy & $\Delta H = -1$           & $\Delta\tau = +1$ \\
Function          & Resolves uncertainty      & Introduces contradiction \\
System effect     & Enables decoding          & Blocks state crystallization \\
Cybernetic analog & Signal in feedback loop   & Unabsorbed disturbance variety \\
Stochastic model  & Binary symmetric channel  & Poisson contradiction process \\
\bottomrule
\end{tabular}
\end{table}

The cybernetic column is significant: in Ashby's framework, a regulator fails when disturbance
variety exceeds its own variety. Torsion is precisely the information-geometric measure of
unabsorbed variety---the contradictions a system has not yet resolved. High torsion is the
collapse-domain equivalent of a regulator that lacks requisite variety: the system cannot
transition to a determinate state because its internal contradictions exceed its resolution
capacity.

The stochastic model row connects to Section~\ref{sec:poisson}: just as Shannon's bit operates
within the binary symmetric channel model with known error probability $p$, the anti-bit operates
within the Poisson contradiction model with known arrival rate $\alpha$ and decay scale $\lambda$.
Both are analytically tractable base cases for their respective domains.

% ============================================================
\section{The Cybernetic Foundation}
% ============================================================

\subsection{Collapse as Generalized Homeostatic Boundary Crossing}

Classical cybernetics describes systems that maintain essential variables within viable limits
through negative feedback. The homeostatic boundary separates viable from non-viable regions of
state space. Ashby's ultrastable system (1960) achieves this by exploring parameter configurations
until essential variables return to the viable zone~\cite{ashby1960}.

The collapse operator generalizes this logic. In homeostasis, the boundary separates viable from
non-viable. In information collapse, the boundary $R - \tau = \varepsilon$ separates latent
(superposed, probabilistic) from collapsed (determinate, irreversible). Both are governed by the
relationship between a system's internal coherence and the threshold it must cross.

\begin{table}[h]
\centering
\caption{Structural correspondence between cybernetic regulation and information collapse.}
\begin{tabular}{>{\raggedright\arraybackslash}p{0.2\textwidth}
                >{\raggedright\arraybackslash}p{0.22\textwidth}
                >{\raggedright\arraybackslash}p{0.45\textwidth}}
\toprule
\textbf{Cybernetic Concept} & \textbf{Collapse Analog} & \textbf{Formal Correspondence} \\
\midrule
Essential variables & State manifold $M$ &
  Latent information space carrying all potential outcomes \\
Negative feedback & Resonance coupling $R$ &
  Amplification of coherent signal toward basis projection \\
Disturbance variety & Torsion $\tau$ &
  Unresolved contradictions opposing collapse \\
Requisite variety (Ashby) & Collapse threshold $\varepsilon$ &
  Minimum net coherence for phase transition \\
Homeostatic boundary & Collapse boundary $R - \tau = \varepsilon$ &
  Hypersurface separating latent from collapsed \\
Ultrastability (Ashby) & Metastability at critical torsion &
  System explores configurations near boundary \\
Circular causality (Wiener) & $R$--$\tau$ co-evolution &
  Resonance and torsion mutually constrain collapse trajectory \\
\bottomrule
\end{tabular}
\end{table}

\subsection{Beyond Maintenance: The Irreversibility Distinction}

The critical distinction between cybernetic regulation and information collapse is
irreversibility. Homeostasis is a maintenance operation: the system oscillates around a setpoint,
and deviations are corrected. The feedback loop is continuous and reversible. Collapse is a phase
transition: the system crosses the boundary $R - \tau \geq \varepsilon$ and the dimensional
reduction from $M$ to $S$ is non-injective (Theorem~\ref{thm:irreversibility}). Information about
the prior superposition is structurally lost.

This distinction maps onto Bateson's hierarchy of learning types~\cite{bateson1972}.
Zero-learning and Learning I are homeostatic---the system maintains or adjusts within a fixed
frame. Learning II (deutero-learning) involves a change of frame itself, which corresponds to
collapse: the system crosses a boundary and its state space is irreversibly restructured. Learning
III---the rare reorganization of the premises of Learning II---corresponds to what we would term a
second-order collapse event, where the collapse operator's own parameters (particularly $B$, the
symbolic basis) are themselves transformed.

\subsection{Wiener's Circular Causality and \texorpdfstring{$R$--$\tau$}{R-tau} Co-evolution}

Wiener's core insight was that cause and effect are not linear but circular~\cite{wiener1948}. In
the collapse operator, resonance $R$ and torsion $\tau$ exhibit precisely this structure.
Increasing resonance with a particular basis element $b$ can resolve contradictions (reducing
$\tau$), while unresolved torsion distorts the resonance landscape (reducing effective $R$). The
trajectory toward or away from collapse is therefore not a simple threshold test but a dynamical
co-evolution:
\begin{align}
  \frac{d\bar\tau}{dt} &\approx \frac{\alpha}{\lambda^2} - \frac{1}{\lambda}\bar\tau
    && \text{(mean-field torsion, Corollary~\ref{cor:meanfield})} \label{eq:coevo-tau} \\
  \frac{dR_{\mathrm{eff}}}{dt} &= \gamma\,(R_{\max} - R_{\mathrm{eff}}) - \eta\,\bar\tau
    && \text{(resonance with torsion drag)} \label{eq:coevo-R}
\end{align}
where $R_{\max}$ is the maximum attainable resonance in the absence of torsion, $\gamma$ is the
resonance growth rate, and $\eta$ is the torsion-drag coefficient quantifying how unresolved
contradictions degrade resonance coupling.

This feedback structure predicts the metastable regime where systems hover near the collapse
boundary without crossing it---the information-geometric equivalent of a cybernetic system in
persistent oscillation around a setpoint it cannot stably reach.

% ============================================================
\section{The Collapse Boundary}
% ============================================================

\subsection{Geometric Interpretation}

The collapse boundary $R - \tau = \varepsilon$ defines a hypersurface in information space
partitioning the state manifold into two regions.

\textbf{Latent region} ($R - \tau < \varepsilon$): high-dimensional dynamics on $M$; superposed
states; multiple potential meanings; probabilistic evolution.

\textbf{Collapsed region} ($R - \tau > \varepsilon$): low-dimensional projection onto $B$; single
determinate state $S$; unique actualized meaning; deterministic outcome.

By Theorem~\ref{thm:irreversibility}, the transition from latent to collapsed is irreversible:
once $P_B$ projects $\psi$ onto $b^*$, the equivalence class $[b^*]$ absorbs all information about
the specific pre-image, and no operation on the collapsed state can recover the full prior
superposition.

\subsection{Phase Structure}

In $(R, \tau)$ space, three regimes are distinguished:
\begin{enumerate}
  \item \textbf{Collapse zone} ($R$ high, $\tau$ low): meaning stabilizes rapidly. The system has
    resolved its contradictions and resonance dominates.
  \item \textbf{Critical region} ($R \approx \varepsilon + \tau$): high sensitivity to
    perturbation; bifurcation behavior. The coupled dynamics
    (Eqs.~\eqref{eq:coevo-tau}--\eqref{eq:coevo-R}) determine whether the system tips toward
    collapse or retreats to latency.
  \item \textbf{Blocked region} ($\tau$ high regardless of $R$): collapse is structurally
    impossible. Torsion variance (Proposition~\ref{prop:moments}) is large, indicating persistent
    instability.
\end{enumerate}

The critical region corresponds to Ashby's ultrastable exploration zone---the system is maximally
sensitive and its outcome is determined by small perturbations in the $R$--$\tau$ co-evolution.

% ============================================================
\section{Substrate Independence}\label{sec:substrate}
% ============================================================

The substrate independence claim asserts that the collapse operator structure
$(M, T, R, B, \varepsilon, \tau) \to S$ and the boundary condition $R - \tau \geq \varepsilon$
appear across fundamentally different domains. We identify the collapse operator structure as
already implicit in several independent lines of published research. In each case, the
six-component structure and the boundary dynamics are present but have not been unified under a
common formalism.

\subsection{Quantum Measurement}\label{sec:quantum}

The quantum measurement problem is the canonical collapse phenomenon. A quantum system in
superposition $|\psi\rangle = \sum_i c_i |b_i\rangle$ collapses to a single eigenstate
$|b_j\rangle$ upon measurement. The mapping to the collapse operator is exact.

\begin{table}[h]
\centering
\caption{Explicit six-component mapping: quantum measurement~\cite{zurek2003} $\to$ collapse
operator.}
\begin{tabular}{>{\raggedright\arraybackslash}p{0.17\textwidth}
                >{\raggedright\arraybackslash}p{0.2\textwidth}
                >{\raggedright\arraybackslash}p{0.5\textwidth}}
\toprule
\textbf{Component} & \textbf{Collapse Operator} & \textbf{Quantum Measurement} \\
\midrule
$M$ & State manifold &
  Hilbert space $\mathcal{H}$ (projective; effectively compact for finite-dim.\ systems) \\
$T$ & Transformation operator & Unitary evolution $U(t) = e^{-iHt/\hbar}$ \\
$R(\psi, b)$ & Resonance coupling &
  Born probability $|c_j|^2 = |\langle b_j|\psi\rangle|^2$ (strictly log-concave in
  Fubini--Study distance) \\
$B$ & Symbolic basis &
  Eigenstates $\{|b_1\rangle, \ldots, |b_k\rangle\}$ of observable $\hat{O}$ \\
$\varepsilon$ & Collapse threshold &
  Decoherence threshold: minimum $|c_j|^2$ for einselection \\
$\tau$ & Torsion &
  Environmental decoherence rate: off-diagonal density-matrix decay $\Gamma_{\mathrm{dec}}$ \\
$S$ & Collapsed state & Einselected pointer state $|b_j\rangle$ \\
\midrule
Boundary & $R - \tau \geq \varepsilon$ &
  $|c_j|^2 - \Gamma_{\mathrm{dec}} \geq \varepsilon_{\mathrm{dec}}$ \\
Irreversibility & Theorem~\ref{thm:irreversibility} &
  Non-invertibility of state reduction $\mathcal{H} \to |b_j\rangle$ \\
Torsion monotonicity & Theorem~\ref{thm:monotonicity} &
  Higher decoherence rates suppress clean collapse \\
\bottomrule
\end{tabular}
\end{table}

Every assumption in Section~2.1 is satisfied: $\mathcal{H}$ (projective) is compact for
finite-dimensional systems; unitary evolution is measure-preserving; the Born rule is strictly
log-concave in Fubini--Study distance; non-degenerate eigenspectra satisfy general position;
decoherence rates are bounded. This is not an analogy---it is a structural isomorphism.

\subsection{Perceptual Bistability}

In perceptual bistability (Necker cube, binocular rivalry), the visual system maintains two
competing interpretations in superposition until one becomes the perceived image. Leopold and
Logothetis (1999) demonstrated that the transition between percepts follows dynamics consistent
with a noisy threshold-crossing process: dominance durations are gamma-distributed, and the
transition is sudden, not gradual~\cite{leopold1999}. The six-component mapping: $M$ is the
perceptual state space (neural population activity); $T$ is neural dynamics (adaptation and mutual
inhibition); $R$ is neural evidence for each interpretation; $B$ is the set of stable percepts;
$\varepsilon$ is the perceptual dominance threshold; $\tau$ is contradictory evidence from the
competing percept. The finding that dominance durations increase when stimulus ambiguity increases
(raising effective $\tau$) is a direct prediction of Theorem~\ref{thm:monotonicity} and is
empirically confirmed~\cite{leopold1999}.

\subsection{Opinion Dynamics and Consensus Formation}

In bounded-confidence opinion dynamics~\cite{deffuant2000, hegselmann2002}, agents update their
opinions only when the distance to a neighbor's opinion falls below a confidence threshold
$\varepsilon$. The six-component mapping: $M$ is the opinion space $[0, 1]^N$; $T$ is the pairwise
update rule; $R$ is the confidence-weighted interaction strength; $B$ is the set of stable
consensus clusters; $\varepsilon$ is the confidence bound; $\tau$ is opinion fragmentation
(variance of the opinion distribution). The sharp transition from plurality to consensus at
critical confidence thresholds is the collapse boundary. Hegselmann and Krause (2002) report that
this transition is abrupt and hysteretic---consistent with the phase-transition character of
$R - \tau = \varepsilon$.

\subsection{Belief Revision and Insight}

In cognitive science, belief revision exhibits threshold dynamics: beliefs resist change under
moderate contradictory evidence but undergo sudden restructuring when evidence accumulates past a
tipping point~\cite{kunda1990, ohlsson2011}. The six-component mapping: $M$ is the belief state
space; $T$ is the evidence accumulation process; $R$ is the fit between evidence and a candidate
belief revision; $B$ is the set of stable belief configurations; $\varepsilon$ is the revision
threshold (the ``impasse barrier'' in Ohlsson's terminology); $\tau$ is contradictory evidence
supporting the current belief (motivated reasoning, in Kunda's framework). The collapse operator
provides a geometric mechanism for what these theories describe phenomenologically: insight occurs
when $R - \tau$ crosses $\varepsilon$, and the experience of suddenness reflects the
phase-transition character of the boundary crossing (Theorem~\ref{thm:irreversibility}).

\subsection{Summary}

The collapse operator structure is not imposed on these domains from outside---it is already
present in each, formalized independently by researchers working within their respective fields.
The quantum measurement mapping is exact at the level of mathematical structure; the perceptual,
social, and cognitive mappings are structural isomorphisms at the level of the six-component
architecture and boundary dynamics, with the specific functional forms of $R$ and $\tau$ varying by
domain. What the Post-Shannon Law contributes is the unifying recognition that these are instances
of a single geometric law.

\begin{definition}[Geometric Coherence Score]
To quantify the structural alignment between a domain instantiation and the universal collapse
operator, we define the GCS. Let $F = (M_d, T_d, R_d, B_d, \varepsilon_d, \tau_d)$ be a
domain-specific instantiation and $\Phi_0 = (M, T, R, B, \varepsilon, \tau)$ the universal
template. The GCS is:
\begin{equation}\label{eq:gcs}
  \mathrm{GCS}(F, \Phi_0) = \frac{1}{6} \sum_{j=1}^{6} s_j(F, \Phi_0),
\end{equation}
where each $s_j \in [0, 1]$ scores the structural correspondence of the $j$-th component.
Independent scoring by domain experts or systematic structural analysis can assess each $s_j$. We
propose this metric for future cross-domain validation studies.
\end{definition}

% ============================================================
\section{Relationship to Prior Work}
% ============================================================

The Post-Shannon Law occupies a specific position in the landscape of information theory
extensions. Shannon (1948) solved transmission~\cite{shannon1948}. Wiener (1948) solved feedback
regulation~\cite{wiener1948}. Ashby (1956) proved the variety constraint on
regulation~\cite{ashby1956}. Bateson (1972) extended cybernetic reasoning to epistemology and
learning~\cite{bateson1972}. Tononi (2004) proposed integrated information ($\Phi$) as a measure of
consciousness~\cite{tononi2004}. Friston (2010) proposed the free-energy principle as a unifying
account of biological self-organization~\cite{friston2010}.

The present framework addresses a gap that none of these fully formalize: the geometric conditions
for irreversible transition from superposed to determinate states. Shannon excludes semantics;
Wiener and Ashby describe maintenance, not phase transition; Bateson identifies the learning
hierarchy but does not formalize the collapse mechanism; Tononi's $\Phi$ measures integration but
not transition dynamics; Friston's free-energy principle describes the minimization drive but not
the boundary conditions for irreversible state crystallization. The collapse operator is
complementary to each of these---it formalizes the boundary crossing that each implies but does
not specify.

% ============================================================
\section{Testable Predictions}
% ============================================================

\subsection{Quantum Measurement}

Systems in the critical torsion regime ($R \approx \varepsilon + \tau$, with $\tau$ scaling with
environmental decoherence) should exhibit classification variance predicted by
Proposition~\ref{prop:moments}: of order $\sqrt{\alpha/3\lambda^3}$ at the stationary
distribution. This predicts measurable fluctuation in einselection outcomes for quantum systems
held near the decoherence threshold, a prediction consistent with existing experimental data on
pointer-state stability under varying environmental coupling~\cite{zurek2003}.

\subsection{Cognitive Science}

The collapse boundary provides a formal account of why meaning formation is experienced as sudden
rather than gradual: insight occurs when the $R$--$\tau$ trajectory crosses the boundary, producing
an irreversible phase transition from ambiguity to clarity. This prediction is testable: perceptual
bistability experiments should show that increasing stimulus ambiguity (raising effective $\tau$)
lengthens dominance durations monotonically, with the functional form matching
Theorem~\ref{thm:monotonicity}. Existing data on Necker cube reversal rates under varying
ambiguity conditions~\cite{leopold1999} provides a candidate dataset.

\subsection{Opinion Dynamics}

The framework predicts that consensus formation in bounded-confidence models should exhibit a
sharp phase transition at the collapse boundary, with the critical confidence threshold
$\varepsilon$ predictable from the torsion structure (opinion fragmentation) of the initial
population. This is testable against existing simulation data from Deffuant et al.\ and
Hegselmann--Krause models~\cite{deffuant2000, hegselmann2002}, and against empirical data on group
deliberation dynamics.

\subsection{AI Safety}

A pre-registered empirical protocol for measuring $R$, $\tau$, and collapse dynamics in human--AI
interaction is specified under CDOT 2.0 (osf.io/7mejx). The protocol operationalizes coupling
coherence, drift pressure, autonomy, and coupling strength as measurable quantities, enabling tests
of Theorems~\ref{thm:irreversibility}--\ref{thm:monotonicity} in the AI-interaction substrate.
Specific predictions: (i) systems in the critical torsion regime should exhibit elevated
classification variance, (ii) uniform torsion shifts should produce monotone decreases in collapse
probability, and (iii) the metastable regime should be detectable as persistent oscillation in
$R$--$\tau$ co-evolution without crossing the boundary.

\subsection{Information Systems Design}

The framework implies design principles: minimize torsion accumulation (reduce internal
contradiction in system states), optimize resonance coupling (ensure system outputs align with
well-defined symbolic bases), and monitor proximity to the collapse boundary (detect when systems
are in the critical region). These translate directly to Ashby's cybernetic design
principle---build systems with sufficient variety to absorb their operating disturbances---but
extend it to the collapse domain: build systems whose torsion resolution capacity matches or
exceeds their contradiction generation rate.

% ============================================================
\section{Open Questions}
% ============================================================

\textbf{Threshold determination.} What governs the value of $\varepsilon$ across substrates? Is
there a universal relationship between manifold dimensionality and collapse threshold, or is
$\varepsilon$ substrate-specific? The Poisson torsion model suggests that $\varepsilon$ may scale
with $\alpha/\lambda$ (the stationary torsion mean), but this requires empirical verification.

\textbf{Empirical torsion measurement.} Can $\tau$ be measured directly in physical systems, or
only inferred from collapse dynamics? The variance bound
$\operatorname{Var}[\tau_\infty] = \alpha/3\lambda^3$ (Proposition~\ref{prop:moments}) provides a
testable prediction: systems with higher contradiction arrival rates should exhibit wider torsion
fluctuations.

\textbf{Quantum mechanical connection.} The structural parallel between the collapse operator and
quantum measurement is explicit but the formal relationship is unresolved. Is the collapse boundary
derivable from quantum decoherence dynamics, or is it an independent geometric principle that
quantum mechanics instantiates?

\textbf{First-principles derivation.} Can the six-component structure be derived from
information-geometric axioms (Fisher metric, Amari $\alpha$-connections) rather than identified
empirically? The strict log-concavity assumption (Assumption~\ref{asm:resonance}) is consistent
with Gaussian kernels on statistical manifolds, suggesting a natural embedding in the Fisher--Rao
framework.

\textbf{Torsion drag coefficient.} The coupled dynamics
(Eqs.~\eqref{eq:coevo-tau}--\eqref{eq:coevo-R}) introduce the parameter $\eta$ (torsion drag on
resonance). Empirical estimation of $\eta$ across domains would test whether the $R$--$\tau$
feedback structure is universal or domain-specific.

% ============================================================
\section{Conclusion}
% ============================================================

Shannon formalized transmission. Cybernetics formalized maintenance. This paper formalizes
collapse: the geometric conditions under which latent information transitions irreversibly into
determinate structure. The Post-Shannon Law---$R - \tau \geq \varepsilon$---specifies the boundary
condition for this phase transition and identifies torsion as the primary structural antagonist of
reality formation.

The framework is grounded in the cybernetic tradition but extends it: where Ashby's requisite
variety governs what a regulator must possess to maintain homeostasis, the collapse threshold
governs what must be exceeded for latent information to become real. The mathematical core rests
on explicit regularity assumptions (compact Riemannian manifold, strictly log-concave resonance,
Poisson contradiction process), yielding provable uniqueness, irreversibility, and torsion
monotonicity. Substrate independence is evidenced by structural isomorphism with quantum
measurement, perceptual bistability, opinion dynamics, and belief revision---each an independent
line of research in which the collapse operator structure is already implicit.

The result is a completion claim: Shannon $+$ cybernetics $+$ collapse $=$ a complete geometric
account of information from creation through transmission through regulation through
crystallization into reality.

\paragraph{Acknowledgments.}
This paper was developed through sustained human-AI cognitive coupling, a methodology central to
the first author's broader research program on Coupling Dynamics (pre-registered at
osf.io/7mejx). The geometric framework, substrate-independence claims, and research direction
originate with the first author. Mathematical formalization, proof construction, and iterative
error-catching --- including correction of the mean-field derivation
(Corollary~\ref{cor:meanfield}) and the torsion monotonicity proof
(Theorem~\ref{thm:monotonicity}) --- were developed through dialogue with the second author
(Claude, Anthropic). Gemini (Google DeepMind) and GPT-4 (OpenAI) contributed earlier-stage
cross-checks on the substrate-independence mappings. The coupling methodology is not incidental to
this contribution and is documented as part of the research program.

\begin{thebibliography}{99}

\bibitem{ashby1956}
Ashby, W.R. (1956). \emph{An Introduction to Cybernetics}. Chapman \& Hall.

\bibitem{ashby1960}
Ashby, W.R. (1960). \emph{Design for a Brain: The Origin of Adaptive Behaviour}, 2nd ed. Chapman
\& Hall.

\bibitem{bateson1972}
Bateson, G. (1972). \emph{Steps to an Ecology of Mind}. University of Chicago Press.

\bibitem{daley2003}
Daley, D.J. \& Vere-Jones, D. (2003). \emph{An Introduction to the Theory of Point Processes},
Vol.~I, 2nd ed. Springer.

\bibitem{deffuant2000}
Deffuant, G., Neau, D., Amblard, F., \& Weisbuch, G. (2000). Mixing beliefs among interacting
agents. \emph{Advances in Complex Systems}, 3(01n04):87--98.

\bibitem{friston2010}
Friston, K. (2010). The free-energy principle: A unified brain theory? \emph{Nature Reviews
Neuroscience}, 11(2):127--138.

\bibitem{hegselmann2002}
Hegselmann, R. \& Krause, U. (2002). Opinion dynamics and bounded confidence: Models, analysis and
simulation. \emph{Journal of Artificial Societies and Social Simulation}, 5(3).

\bibitem{kingman1993}
Kingman, J.F.C. (1993). \emph{Poisson Processes}. Oxford University Press.

\bibitem{kunda1990}
Kunda, Z. (1990). The case for motivated reasoning. \emph{Psychological Bulletin},
108(3):480--498.

\bibitem{leopold1999}
Leopold, D.A. \& Logothetis, N.K. (1999). Multistable phenomena: Changing views in perception.
\emph{Trends in Cognitive Sciences}, 3(7):254--264.

\bibitem{ohlsson2011}
Ohlsson, S. (2011). \emph{Deep Learning: How the Mind Overrides Experience}. Cambridge University
Press.

\bibitem{petersen2016}
Petersen, P. (2016). \emph{Riemannian Geometry}, 3rd ed. Graduate Texts in Mathematics, vol.~171.
Springer.

\bibitem{shannon1948}
Shannon, C.E. (1948). A mathematical theory of communication. \emph{Bell System Technical
Journal}, 27(3):379--423.

\bibitem{tononi2004}
Tononi, G. (2004). An information integration theory of consciousness. \emph{BMC Neuroscience},
5(42).

\bibitem{wiener1948}
Wiener, N. (1948). \emph{Cybernetics: Or Control and Communication in the Animal and the Machine}.
MIT Press.

\bibitem{zurek2003}
Zurek, W.H. (2003). Decoherence, einselection, and the quantum origins of the classical.
\emph{Reviews of Modern Physics}, 75(3):715--775.

\end{thebibliography}

\end{document}
